\documentclass{article}
\usepackage{graphicx}
\usepackage{amsmath,amssymb}
\usepackage{booktabs}
\usepackage{hyperref}

\title{Preserving Strategy Specialization in Shared Multi-Robot Policies}
\author{TheK-B1000}
\date{September 2026}

\begin{document}

\maketitle

%======================================================================
\section*{Abstract}
%======================================================================

Multi-robot teams operating against heterogeneous adversaries may require
qualitatively different coordination strategies under different opponent
regimes. Deploying a separate network for every regime preserves specialization
but scales poorly; a single shared policy conditioned on a discrete strategy
code is more compact, yet strategically important behavior can be difficult not
only to learn, but also to preserve under continued optimization and parameter
sharing. We study strategy preservation in a partially observable maritime
multi-robot Capture-the-Flag task with two scripted opponent regimes. We
transfer complementary expert strategies into strategy-conditioned policies and
systematically increase parameter sharing from perception to higher-level action
components. We additionally introduce \emph{Strategy-Preserving PPO} (SP-PPO),
combining \emph{Counterfactual Strategy Correction} (CSC) with
\emph{Strategy-Preserving Fine-Tuning} (SPFT), as an intervention family for
studying whether strategically consequential decisions can be identified and
retained during continued optimization. Evaluation separates action imitation,
policy separation, and payoff-level strategic specialization. In the
progressive-sharing study, sharing only the perception encoder preserves the
specialization criterion with no detectable paired loss relative to independent
experts, while deeper sharing yields detectable or fit-confounded degradation.
These results clarify which components of a multi-robot policy can be shared
while attempting to retain opponent-conditioned coordination strategies under
the stated training and evaluation protocol.

%======================================================================
\section{Introduction}
%======================================================================

Multi-robot systems that operate against heterogeneous environments or
adversaries often cannot rely on a single coordination style. Different opponent
behaviors can require qualitatively different team strategies---for example, a
broad zone defense may favor one form of team organization, while a compact,
reactive defense favors another. Learning policies that express the right
strategy for the encountered regime is therefore a central robotics problem,
not merely a representation-learning detail.

Storing a separate neural policy for every operating regime preserves
specialization but does not scale well as the number of regimes grows. Memory,
deployment complexity, and maintenance cost all increase with the number of
independent networks. A natural alternative is one shared multi-robot policy
conditioned on a small discrete \emph{strategy code}: the code selects which
opponent-conditioned behavior the shared policy should express.

Parameter sharing, however, can erase strategically important differences even
when average action imitation remains high. Closed-loop coordination depends on
temporally extended decisions; agreement on individual actions does not
automatically imply that each conditioned mode retains its intended task
advantage. The scientific question of this paper is therefore:

\begin{quote}
\emph{Can a single strategy-conditioned multi-robot policy preserve multiple
opponent-specific strategies, and how much neural-network parameter sharing can
be introduced before those strategies interfere with one another?}
\end{quote}

Existing latent-policy and multi-agent reinforcement learning methods have shown
that a shared policy can represent diverse behaviors, roles, or
opponent-conditioned responses. However, substantially less attention has been
paid to whether strategically distinct behaviors remain preserved as
increasingly large portions of the policy are forced to share parameters. This
distinction is important because high action-level imitation fidelity does not
necessarily imply preservation of closed-loop strategic behavior.

We therefore study strategy preservation under progressive policy sharing in a
multi-robot adversarial task. We first establish independently trained expert
policies with complementary opponent-conditioned behavior, then transfer these
strategies into strategy-conditioned policies while systematically increasing
parameter sharing from perception to higher-level action components. In
addition, we introduce \emph{Strategy-Preserving PPO} (SP-PPO) as a
methodological intervention family:
\emph{Counterfactual Strategy Correction} (CSC) identifies which strategically
important decisions should be corrected using matched counterfactual rollouts,
and \emph{Strategy-Preserving Fine-Tuning} (SPFT) aims to prevent those
corrections from being erased during continued PPO optimization. Our evaluation
separates action imitation, policy separation, and payoff-level strategic
specialization, allowing us to identify when a policy remains locally similar to
its experts while losing behaviorally meaningful strategy specialization.

The main contributions of this work are:
\begin{enumerate}
  \item We introduce a progressive sharing study that isolates the effect of
  sharing perception, policy-backbone, and action-output parameters on
  opponent-conditioned strategy specialization while holding expert
  supervision, training data, optimization budget, and evaluation conditions
  fixed.
  \item We distinguish action-level imitation fidelity from closed-loop
  strategic fidelity and evaluate specialization through matched-seed payoff
  contrasts rather than latent separation or action agreement alone.
  \item We develop \emph{Strategy-Preserving PPO} (SP-PPO), a multi-robot
  policy-learning framework for studying whether counterfactual strategy
  correction and retention-based fine-tuning can preserve opponent-conditioned
  strategies during continued optimization.
  \item Within SP-PPO, we develop \emph{Counterfactual Strategy Correction}
  (CSC), which uses matched continuation rollouts to identify strategically
  consequential action interventions and incorporates positively validated
  interventions as strategy-conditioned corrective supervision.
  \item Within SP-PPO, we develop \emph{Strategy-Preserving Fine-Tuning}
  (SPFT), which augments CSC with temporal policy retention to reduce the
  washout of sparse strategy-defining updates during continued PPO
  optimization.
\end{enumerate}

%======================================================================
\section{Problem Setup}
\label{sec:setup}
%======================================================================

\subsection{Environment and agents}

We study a partially observable, adversarial multi-robot setting instantiated
as a nominal $2$-vs-$2$ naval Capture-the-Flag task on a symmetric maritime
arena. Blue is the learned multi-robot team; Red is a scripted or frozen
adversarial opponent. Each side controls $|B|=|R|=2$ autonomous surface vessels.
We model the interaction as a finite-horizon partially observable Markov game
\[
G =
\left(
\mathcal{N},\mathcal{S},\{\mathcal{A}_i\}_{i\in\mathcal{N}},P,
\{r_i\}_{i\in\mathcal{N}},\{\Omega_i\}_{i\in\mathcal{N}},\{O_i\}_{i\in\mathcal{N}},\gamma
\right),
\]
with $\mathcal{N}=B\cup R$.

At each simulator tick $t$, agent $i$ receives a local observation
$o_{i,t}\sim O_i(\cdot\mid s_t)$ and may request a temporally extended
\emph{macro-action} together with a target. The macro vocabulary is
$\{\texttt{GO\_TO},\texttt{GRAB\_MINE},\texttt{GET\_FLAG},\texttt{PLACE\_MINE},\texttt{GO\_HOME}\}$.
For reproducibility, the joint action is represented as
$\mathrm{MultiDiscrete}([5,50,5,50])$ (per agent: a $5$-way macro and a
$50$-way target). The joint action induces $s_{t+1}\sim P(\cdot\mid s_t,a_t)$.

Each team's flag is anchored at its home base; a carrier that returns the enemy
flag alive scores a capture. Mines are stationary hazards that create
area-denial and ambush dynamics.

\subsection{Asynchronous macro-action decision points}

An accepted macro-action is held for a macro-specific number of ticks before the
agent may choose a new one
($\texttt{GO\_TO}=4$, $\texttt{GRAB\_MINE}=3$, $\texttt{GET\_FLAG}=4$,
$\texttt{PLACE\_MINE}=2$, $\texttt{GO\_HOME}=4$). Agent $i$ carries a countdown
$c_{i,t}$; a newly requested macro is accepted only when
\[
d_t^i=\mathbb{I}\!\left[c_{i,t}\le 0\right].
\]
When $d_t^i=0$, the legal-action mask collapses to the already-committed macro,
so the network output is not an executable decision. Because hold lengths differ
by macro and agents commit independently, decision boundaries are asynchronous.
For policy evaluation and supervision, the task is therefore better viewed as a
semi-Markov decision process over commitment decisions than as a flat per-tick
Markov decision process.

\subsection{Opponent regimes}

Red is drawn from one of two scripted defensive philosophies that we call
\emph{opponent regimes}:
\begin{itemize}
  \item \textbf{Opponent Regime~A:} broad / passive zone defense.
  \item \textbf{Opponent Regime~B:} compact / reactive defense.
\end{itemize}
The two regimes are designed to elicit contrasting strategic responses from
Blue rather than represent two difficulty levels of the same opponent behavior.
Implementation identifiers used for reproducibility (configuration overlay
\textbf{OP6} for Regime~A; canonical \textbf{OP7} for Regime~B), together with
frozen behavior-tree parameters such as
$\texttt{defender\_zone\_frac}$ and $\texttt{threat\_radius}$, are recorded in
the experiment configuration and are not repeated as conceptual labels in the
main narrative.

We do not assume a priori that the two regimes require different Blue
strategies. An \emph{opponent-conditioned strategy-preference test} measures
whether a GUARD-biased response is preferred over a BREACH-biased response under
Regime~A, and the reverse contrast under Regime~B, with pass condition
$\mathrm{LCB}_{95}>0$ on each contrast. This empirical test establishes whether
the regimes induce complementary strategic demand. If one Blue policy were
simply better against both regimes, a single shared behavior would be
appropriate rather than a specialization problem.

\subsection{Strategy-conditioned policies}

Let $\pi_\theta(a\mid o,z)$ denote a shared multi-robot policy conditioned on a
discrete \emph{strategy code} $z\in\{z_0,z_1\}$. The strategy code selects which
opponent-conditioned behavior the shared policy should express. The intended
deployment mapping is $z_0\mapsto\text{Regime~A}$ and $z_1\mapsto\text{Regime~B}$;
once fixed, this assignment is never flipped. Let $V(z,p)$ denote the win rate
when deploying with strategy code $z$ against opponent regime
$p\in\{A,B\}$, estimated over a held-out seed block.

Independently trained \emph{strategy-specific expert policies} $\pi_A$ and
$\pi_B$ provide the zero-sharing reference: each is trained against one regime
with no shared parameters and no strategy-code mechanism. A single
\emph{generalist} policy $\pi_G$ trained on the mixed opponent distribution
without a strategy code provides a single-policy baseline.

\subsection{Strategy-specialization criterion}
\label{sec:specialization}

We require more than behavioral diversity. For two strategy codes to represent
meaningful specialization, each code must provide higher task success than the
alternative code in its intended opponent regime. A policy is considered
strategy-specialized only if each strategy code significantly outperforms the
alternative code in the opponent regime for which it is intended.

Formally, define
\[
\Delta_A = V(z_0,A)-V(z_1,A),
\qquad
\Delta_B = V(z_1,B)-V(z_0,B).
\]
The \emph{strategy-specialization criterion} (equivalently, the two-sided
\emph{specialization crossover} pattern) requires
\[
\mathrm{LCB}_{95}(\Delta_A)>0
\quad\wedge\quad
\mathrm{LCB}_{95}(\Delta_B)>0,
\]
evaluated by a paired percentile bootstrap over the held-out seeds with
$n_{\mathrm{boot}}=20{,}000$, $\alpha=0.05$, and the seed as the resampling
unit. The binary win indicator is the primary outcome because it is the
deployment quantity used by every specialization evaluation.

Independent-expert specialization is reported on a held-out $64$-seed block.
All parameter-sharing comparisons use the matched $128$-seed paired protocol of
Section~\ref{sec:matched}.

\subsection{Scope of claims}

Primary claims apply to the nominal $2$-vs-$2$ environment, the two opponent
regimes above, the specific training recipe, and the fixed evaluation
implementation. Scaling to $4$-vs-$4$ and $6$-vs-$6$, and robustness
perturbations such as localization noise, motion uncertainty, and control delay,
are external-validity experiments. Until completed, they are not used as evidence
for the primary sharing conclusions and are not used to select the primary
architecture.

%======================================================================
\section{Methodology}
\label{sec:method}
%======================================================================

\subsection{Reference policies}

\textbf{Independent experts (zero-sharing reference).}
We train strategy-specific experts $\pi_A$ and $\pi_B$ under matched conditions,
one per opponent regime, with no shared parameters and no strategy-code
mechanism. The independently trained experts provide the zero-sharing reference
for strategy specialization.

\textbf{Generalist baseline.}
We also evaluate $\pi_G$, a single policy trained on the mixed opponent
distribution with no strategy-code mechanism. The generalist provides a
single-policy baseline trained on the mixed opponent distribution.

\subsection{Strategy-conditioned policy distillation}
\label{sec:distill}

We transfer expert behaviors into a shared strategy-conditioned policy
$\pi_\theta(a\mid o,z)$ by supervised distillation. Both experts are queried on
the \emph{same} stored states, so differences between conditioned outputs are
attributable to expert strategies rather than different state datasets.

The distillation objective is
\[
\mathcal{L}_{\mathrm{distill}}
=
\tfrac{1}{2}D_{\mathrm{KL}}\!\left(\pi_A(\cdot\mid s)\,\|\,\pi_\theta(\cdot\mid s,z_0)\right)
+
\tfrac{1}{2}D_{\mathrm{KL}}\!\left(\pi_B(\cdot\mid s)\,\|\,\pi_\theta(\cdot\mid s,z_1)\right),
\]
evaluated per action head over masked distributions and averaged over
decision-eligible heads only ($d_t^i=1$). Only actor parameters are optimized;
critic parameters remain unchanged and are verified immobile. Training runs a
fixed number of epochs with no early stopping and no checkpoint selection; the
terminal checkpoint is the sole artifact.

We maintain three measurement levels and do not conflate them:
\begin{itemize}
  \item \textbf{Action-level imitation:} argmax agreement and KL divergence
  relative to the matched expert.
  \item \textbf{Strategy separation:} Jensen--Shannon divergence between the
  two conditioned action distributions.
  \item \textbf{Closed-loop strategy specialization:} $V(z,p)$, $\Delta_A$,
  $\Delta_B$, and the paired-bootstrap criterion of
  Section~\ref{sec:specialization}.
\end{itemize}
High action-level imitation does not logically imply task-level specialization.
Different strategy-conditioned action distributions likewise do not logically
imply payoff-relevant specialization. The payoff-based criterion is the primary
task-level definition.

\subsection{Strategy-Preserving PPO}
\label{sec:spppo}

We introduce \emph{Strategy-Preserving PPO} (SP-PPO) as a methodological
intervention family for studying strategy installation and retention in a
shared multi-robot policy. CSC identifies which strategically important
decisions should be corrected using matched counterfactual rollouts. SPFT aims
to prevent those corrections from being erased during continued PPO
optimization.

\subsection{Counterfactual Strategy Correction (CSC)}
\label{sec:csc}

CSC addresses \emph{which decisions should change}. Because macro-actions are
temporally extended, we evaluate interventions only at decision-eligible ticks
($d_t^i=1$). At such boundaries, matched continuation rollouts estimate whether
substituting an expert decision improves terminal task success relative to a
baseline continuation of the current strategy-conditioned policy.

\paragraph{Matched-continuation intervention estimate.}
For a boundary state $s_t$ drawn from a prospectively selected, outcome-blind,
stratified sample of the baseline policy's deployment distribution, and an
estimand $e$ identifying which agent(s) may be intervened on, define three
matched continuation arms that share an identical continuation seed $r_j$:
$R_0$ (baseline continues unmodified), $R_A$ ($\pi_A$ takes over the intervened
agent(s) for the remainder of the episode), and $R_B$ ($\pi_B$ takes over
identically). With $M=16$ matched continuations,
\[
\hat{A}_A(x)=\frac{1}{M}\sum_{j=1}^{M}\left(R_A^j-R_0^j\right),
\qquad
\hat{A}_B(x)=\frac{1}{M}\sum_{j=1}^{M}\left(R_B^j-R_0^j\right),
\]
where $R$ is the binary episode win indicator. These are matched-continuation
intervention estimates under the controlled simulator experiment, not general
causal effects beyond that design.

\paragraph{Positive-advantage routing.}
An expert is selected only when its estimated advantage is strictly positive
over both zero and the alternative expert:
\[
t^*(x)=
\begin{cases}
\pi_A, & \hat{A}_A(x)>\max\left(0,\hat{A}_B(x)\right),\\[3pt]
\pi_B, & \hat{A}_B(x)>\max\left(0,\hat{A}_A(x)\right),\\[3pt]
\mathrm{none}, & \text{otherwise},
\end{cases}
\qquad
w(x)=
\begin{cases}
\hat{A}_{t^*}(x), & t^*(x)\neq\mathrm{none},\\[3pt]
0, & \text{otherwise}.
\end{cases}
\]
Exact positive ties receive zero weight. At both-free states, a usable joint
intervention supersedes constituent per-agent units.

\paragraph{Counterfactual correction loss.}
The $(o,a)$ trajectory used for supervision is obtained by rolling $t^*(x)$
exactly once under a canonical seed independent of the measurement seeds
$\{r_j\}$, so trajectory construction cannot alter $t^*$ or $w$. The auxiliary
loss is a weighted negative log-likelihood of the routed expert's action,
restricted to decision-eligible agents:
\[
\mathcal{L}_{\mathrm{csc}}
=
\frac{\sum_{t,i} d_t^i\,w_{t,i}\,\mathcal{L}_{\mathrm{expert}}(s_t,i,z)}
{\sum_{t,i} d_t^i\,w_{t,i}}.
\]
Committed agents ($d_t^i=0$) contribute to neither numerator nor denominator.
CSC fine-tuning combines this term with task PPO,
\[
\mathcal{L}=\mathcal{L}_{\mathrm{PPO}}+\lambda_c\mathcal{L}_{\mathrm{csc}},
\]
with $\lambda_c=0.05$. Reward, returns, GAE, and value targets are unchanged.
A matched control arm uses the same warm start, horizon, and PPO recipe without
$\mathcal{L}_{\mathrm{csc}}$.

\subsection{Strategy-Preserving Fine-Tuning (SPFT)}
\label{sec:spft}

SPFT addresses \emph{how sparse corrections are retained during continued
optimization}. CSC interventions are sparse: only decision points with strictly
positive matched-continuation advantage receive nonzero weight. Continued PPO
updates can therefore overwrite strategically consequential corrections. SPFT
keeps the CSC correction bank and $\lambda_c$ fixed and adds temporal retention
against an exponential-moving-average (EMA) actor teacher.

\paragraph{EMA teacher and retention loss.}
Initialize $\bar\theta_0=\theta_0$ as an exact copy of the warm-start actor.
After each actor optimizer update,
\[
\bar\theta\leftarrow 0.995\,\bar\theta + 0.005\,\theta.
\]
The retention loss is a KL between the EMA teacher and the current policy,
\[
\mathcal{L}_{\mathrm{ret}}
=
D_{\mathrm{KL}}\!\left(\pi_{\bar\theta}(\cdot\mid o,z)\,\|\,\pi_\theta(\cdot\mid o,z)\right),
\]
applied only at decision-eligible agents under the same legality masks used by
PPO. Mid-hold (committed) frames contribute zero retention loss.

\paragraph{Full SP-PPO objective.}
When CSC and SPFT are combined, the actor objective is
\[
\mathcal{L}
=
\mathcal{L}_{\mathrm{PPO}}
+
\lambda_c\mathcal{L}_{\mathrm{csc}}
+
\lambda_r\mathcal{L}_{\mathrm{ret}},
\]
with $\lambda_c=0.05$ and $\lambda_r=0.01$. A matched control arm uses CSC
without retention, isolating the SPFT term.

\subsection{Progressive parameter-sharing study}
\label{sec:sharing}

Holding the distillation recipe, dataset, schedule, and evaluation seeds fixed,
we vary only which modules are tied between the two strategy branches. Sharing
is cumulative. The manipulated variable is the location and amount of parameter
sharing.

\begin{table}[t]
\centering
\caption{Progressive parameter-sharing conditions. Sharing is cumulative.
Every condition uses the same experts, distillation data, optimization budget,
evaluation seeds, and specialization criterion.}
\label{tab:sharing}
\begin{tabular}{lll}
\toprule
\textbf{Condition} & \textbf{Shared components} & \textbf{Strategy-specific components} \\
\midrule
Share-0: Independent Experts
  & None
  & Entire actor \\
Share-Encoder: Shared Perception
  & Perception encoder
  & Backbone + outputs \\
Share-Backbone: Shared Perception + Backbone
  & Encoder + MLP backbone
  & Strategy branches + outputs \\
Share-Macro: Shared Backbone + Macro-Action Output
  & Encoder + backbone + macro-action outputs
  & Target-selection outputs \\
\bottomrule
\end{tabular}
\end{table}

Each condition uses complete actors of the experts' own architecture class,
thereby reducing the likelihood that performance differences are explained solely
by reduced per-strategy capacity. Module tying is verified by object identity
and unique-parameter arithmetic. A mixed-$z$ minibatch gives each shared module
gradient from both strategies in the same step.

\paragraph{Share-0.}
Share-0 installs $\pi_A$ and $\pi_B$ behind a strategy-code dispatch wrapper with
zero shared parameters, verified bit-exact against both experts on the
evaluation device. It establishes that the dispatch-and-evaluate path can
express specialization when no parameters are shared.

\subsection{Matched-seed paired evaluation}
\label{sec:matched}

All parameter-sharing conditions are evaluated on the same $128$ held-out
episode seeds as the Share-0 reference. For each seed, we compute the change in
specialization contrast relative to Share-0, yielding a paired estimate that
controls for variation in episode difficulty. The Share-0 outcomes on this
matched set are fixed before evaluating subsequent sharing conditions.

The primary statistics are
\[
D_A=\mathbb{E}_s\!\left[\Delta_A^{\mathrm{(cond)}}(s)-\Delta_A^{\mathrm{(Share\textrm{-}0)}}(s)\right],
\qquad
D_B=\mathbb{E}_s\!\left[\Delta_B^{\mathrm{(cond)}}(s)-\Delta_B^{\mathrm{(Share\textrm{-}0)}}(s)\right],
\]
under a paired bootstrap. We declare a detectable sharing-associated reduction
in specialization on a regime only if $\mathrm{UCB}_{95}(D_p)<0$: that is, the
upper $95\%$ confidence bound on the paired change lies below zero.

Because CPU and CUDA execution can produce different closed-loop trajectories
despite identical nominal seeds, all primary paired comparisons use a fixed CUDA
execution path. Diagnostics collected on CPU are scoped to the CPU-generated
state distribution.

\subsection{Imitation-fidelity control}
\label{sec:fidelity}

To interpret a reduction in specialization as associated with representation
sharing, compared models must reproduce their expert policies with similar
fidelity. We therefore define an imitation-fidelity comparability band:
final holdout argmax agreement must lie within $0.010$ of the Share-Encoder
reference agreements on the same holdout set. This band is an
\emph{interpretation criterion}, not a success gate and not a theoretically
privileged threshold. Continuous agreement and KL values are reported for every
condition regardless of band membership. If a condition falls outside the band,
any specialization degradation is reported as sharing-associated but is not
uniquely attributed to strategy homogenization or specialization loss /
negative transfer. Observed Share-Encoder reference agreements are reported in
Results.

\subsection{Mechanism decomposition}
\label{sec:decomp}

For fidelity-comparable conditions, we decompose specialization change into:
\begin{itemize}
  \item \textbf{On-regime} performance: strategy code evaluated against its
  intended opponent regime.
  \item \textbf{Cross-regime} performance: the same strategy code evaluated
  against the other opponent regime.
\end{itemize}
For fidelity-comparable conditions, increased cross-regime performance with
preserved on-regime performance is interpreted as evidence consistent with
strategy homogenization, whereas decreased on-regime performance is interpreted
as specialization loss / negative transfer. If imitation fidelity is not
comparable, we report sharing-associated degradation without assigning a
mechanism.

\subsection{Mechanism diagnostics}
\label{sec:diagnostics}

To interpret cases in which a distilled strategy-conditioned policy does not
satisfy the specialization criterion despite high imitation fidelity, we apply
three diagnostic analyses:

\paragraph{H1: Distribution shift.}
Does expert--student agreement decrease on student-visited states relative to
expert-visited states?

\paragraph{H2: High-leverage mismatches.}
Are expert--student disagreements disproportionately associated with final
outcomes? Using matched counterfactual branching---replaying a recorded prefix,
seeding the environment identically across branches, and altering exactly one
decision---we compare expert substitution at disagreement states against a
decision-sensitivity reference.

\paragraph{H3: Action-component localization.}
Are residual mismatches concentrated in macro-action selection or in target
selection?

Each diagnostic is a measurement with a fixed interpretation rule, not a
training method. Device scope is stated with each result.

\subsection{Reproducibility}
\label{sec:repro}

Training and evaluation specifications are fixed before execution. Final
comparisons use held-out evaluation seeds and terminal checkpoints without
post-hoc checkpoint selection. Architecture sharing is verified
programmatically, and checkpoint hashes identify evaluated models. Because
execution device can affect closed-loop trajectories, all primary paired
evaluations use a fixed CUDA path.

%======================================================================
\section{Results}
\label{sec:results}
%======================================================================

\subsection{Do the opponent regimes require different strategies?}

The opponent-conditioned strategy-preference test establishes complementary
strategic demand under Regimes~A and~B.
% [draft] Insert quantitative table/figure (GUARD vs BREACH contrasts; LCB95).

\subsection{Can independent expert policies specialize?}

Independent experts $\pi_A$ and $\pi_B$, scored under a preregistered protocol
on a held-out $64$-seed block, satisfy the strategy-specialization criterion:
\[
\Delta_A = +0.172\;[+0.016,+0.328],
\qquad
\Delta_B = +0.297\;[+0.125,+0.469].
\]
The generalist attains
$V(\pi_G,A)=0.375\;[0.265,0.500]$ and
$V(\pi_G,B)=0.500\;[0.375,0.625]$.

On the matched $128$-seed Share-0 reference,
\[
\Delta_A = +0.289\;[+0.164,+0.406],
\qquad
\Delta_B = +0.258\;[+0.141,+0.375].
\]

\subsection{Can those strategies be represented in one conditioned policy?}

\paragraph{Strategy-conditioned distillation.}
The distilled shared policy reproduces its experts closely (holdout argmax
agreement $0.965$ for $z_0\!\sim\!\pi_A$ and $0.984$ for $z_1\!\sim\!\pi_B$;
conditioned-branch Jensen--Shannon divergence $0.183$, compared with the
experts' $0.179$) yet does not satisfy the specialization criterion:
\[
\Delta_A=+0.094\;[-0.078,+0.266],
\qquad
\Delta_B=+0.141\;[-0.047,+0.313].
\]
Both point estimates are correctly signed; neither confidence interval clears
zero. High-fidelity imitation therefore does not automatically preserve
closed-loop specialization on the student's own state distribution.

\paragraph{Mechanism diagnostics.}
On $24{,}693$ student-generated states, expert agreement remains high
($0.971$ / $0.986$), not below expert-visited holdout agreement
($0.965$ / $0.984$), so H1 is not supported in the tested CPU regime.
For H2, the powered $z_0$ arm shows $18$ of $20$ expert substitutions
producing an identical final outcome (mean $|L_{\mathrm{margin}}|=0.100$ vs.\
decision-sensitivity reference $0.450$); the $z_1$ arm ($6$ branch points) is
underpowered and yields no conclusion. For H3, macro disagreement is
$1.11\%$ ($z_0$) and $0.23\%$ ($z_1$), while target disagreement is $5.21\%$
and $1.60\%$ (factors of approximately $4.7$--$7.0$).

\subsection{Strategy-Preserving PPO components}

CSC and SPFT are evaluated under the same strategy-specialization criterion on
held-out seed blocks, with matched treatment/control designs.

\paragraph{Counterfactual Strategy Correction (CSC).}
Neither the CSC treatment arm nor the matched PPO-only control satisfies the
specialization criterion. Point estimates are negative on both regimes for both
arms (treatment:
$\Delta_A=-0.141\;[-0.281,0.000]$,
$\Delta_B=-0.031\;[-0.203,+0.141]$;
control:
$\Delta_A=-0.188\;[-0.328,-0.047]$,
$\Delta_B=-0.063\;[-0.250,+0.125]$).
Paired treatment-minus-control contrasts do not clear zero on either regime.
Thus CSC does not recover payoff-level specialization under the frozen protocol.

\paragraph{Strategy-Preserving Fine-Tuning (SPFT).}
Neither the SPFT treatment nor the CSC-only control satisfies the
specialization criterion (treatment:
$\Delta_A=-0.156\;[-0.313,+0.016]$,
$\Delta_B=+0.031\;[-0.141,+0.203]$;
control:
$\Delta_A=-0.078\;[-0.219,+0.063]$,
$\Delta_B=+0.063\;[-0.109,+0.234]$).
Paired SPFT-minus-control contrasts likewise do not clear zero. Under the stated
protocol, retention stabilization does not establish recovery of specialization
relative to CSC alone.

\subsection{Where does parameter sharing begin to affect specialization?}

\begin{table}[t]
\centering
\caption{Progressive parameter-sharing results on the matched $128$-seed set.
$D_A$ and $D_B$ are within-seed differences against Share-0.}
\label{tab:sharing-results}
\begin{tabular}{lllll}
\toprule
\textbf{Condition} & \textbf{Shared} & $\boldsymbol{\Delta_A}$ & $\boldsymbol{\Delta_B}$ &
\textbf{Paired $D$} \\
\midrule
Share-0
  & none
  & $+0.289\;[+0.164,+0.406]$
  & $+0.258\;[+0.141,+0.375]$
  & reference \\
Share-Encoder
  & CNN encoder
  & $+0.266\;[+0.148,+0.383]$
  & $+0.164\;[+0.047,+0.281]$
  & $D_A=-0.023$, $D_B=-0.094$; no detectable loss \\
Share-Backbone
  & $+$ MLP backbone
  & $+0.148\;[+0.023,+0.266]$
  & $+0.141\;[+0.016,+0.266]$
  & $D_A=-0.141$ (detectable loss);
    $D_B=-0.117$ (no detectable loss) \\
Share-Macro
  & $+$ macro-action outputs
  & $+0.125\;[0.000,+0.250]$
  & $+0.156\;[+0.031,+0.281]$
  & $D_A=-0.164$ (detectable loss);
    $D_B=-0.102$ (no detectable loss) \\
\bottomrule
\end{tabular}
\end{table}

\paragraph{Share-Encoder.}
Share-Encoder holdout argmax agreements are $0.967/0.989$; these values define
the imitation-fidelity reference band used in Section~\ref{sec:fidelity}.
Sharing only the perception encoder preserves the specialization criterion on
both regimes, and neither within-seed difference against Share-0 reaches
$\mathrm{UCB}_{95}(D_p)<0$. Decomposing by role: on-regime cells shift by
$-0.012$ on average, while cross-regime cells shift by $+0.047$, driven by
$z_0$ under Regime~B rising $+0.086$. This pattern is consistent with strategy
homogenization rather than on-regime specialization loss.

\paragraph{Share-Backbone.}
Share-Backbone terminal holdout agreement is $0.963/0.985$ (within $0.010$ of
the Share-Encoder reference on both codes), so fidelity is comparable.
Branch Jensen--Shannon divergence is $0.172$ (Share-Encoder $0.178$; experts
$0.179$). The specialization criterion still holds on both regimes, but the
paired contrast against Share-0 shows a detectable reduction on Regime~A
($\mathrm{UCB}_{95}(D_A)<0$). On-regime performance declines on average while
cross-regime performance rises, a pattern consistent with specialization loss /
negative transfer under fidelity-comparable sharing of the backbone.

\paragraph{Share-Macro.}
Share-Macro terminal holdout agreement is $0.936/0.969$, outside the
imitation-fidelity band anchored on Share-Encoder. The specialization criterion
is not satisfied ($\mathrm{LCB}_{95}(\Delta_A)=0$). Paired evaluation likewise
shows a detectable reduction on Regime~A. Because imitation fidelity is not
comparable, we report this as sharing-associated and fit-confounded degradation
and do not uniquely attribute it to strategy homogenization or specialization
loss.

\subsection{External validity}
\label{sec:external}

$4$-vs-$4$ / $6$-vs-$6$ scalability and robustness perturbations (localization,
motion, control delay) are reported only when completed, and only as
external-validity experiments---not as evidence used to choose among sharing
conditions.

%======================================================================
\section{Discussion}
\label{sec:discussion}
%======================================================================

The primary scientific finding concerns progressive parameter sharing:
perception-level sharing (Share-Encoder) preserves the specialization criterion
under the matched-seed protocol, with a decomposition consistent with mild
strategy homogenization rather than on-regime specialization loss. Deeper
sharing of the MLP backbone (Share-Backbone) remains fidelity-comparable and
still clears the absolute specialization criterion, yet produces a detectable
paired reduction relative to Share-0 on Regime~A, consistent with specialization
loss under additional representation sharing. Further sharing at macro-action
outputs (Share-Macro) coincides with deteriorated imitation fidelity and failure
of the specialization criterion; under the imitation-fidelity control, that
degradation is reported as sharing-associated but not uniquely attributed to a
single interference mechanism.

SP-PPO is a methodological contribution and an informative intervention family
for strategy installation and retention. CSC identifies high-leverage strategic
decisions via matched counterfactual rollouts; SPFT aims to keep those
corrections from drifting away under continued PPO. Under the frozen protocols
studied here, neither CSC nor SPFT satisfies the payoff-based specialization
criterion. Identifying positively advantageous interventions---and attempting to
retain them---therefore does not automatically restore closed-loop complementary
strategies. Local policy similarity remains insufficient evidence of preserved
strategic identity.

For multi-robot policy design, the sharing study suggests that substantial
perception sharing can be tolerable in this setting, while sharing deeper
decision-related representations is where opponent-conditioned specialization
becomes fragile. Claims remain scoped to the two-regime nominal $2$-vs-$2$ task
and training recipes studied here. External validity for larger teams and
sensing/control perturbations remains outside the primary claim until those
experiments are finished.

%======================================================================
\section{Conclusion}
\label{sec:conclusion}
%======================================================================

We study whether a strategy-conditioned multi-robot policy can preserve
opponent-conditioned coordination strategies under progressive parameter sharing
and continued optimization. Independent experts provide a zero-sharing
specialization reference. The progressive-sharing study is the primary empirical
contribution: Share-Encoder preserves specialization, Share-Backbone introduces
a detectable paired specialization loss on one regime despite comparable
fidelity, and Share-Macro yields fit-confounded degradation together with
failure of the specialization criterion. Strategy-Preserving PPO (SP-PPO)
combines Counterfactual Strategy Correction (CSC) with Strategy-Preserving
Fine-Tuning (SPFT) as a methodological intervention family for strategy
installation and retention; under the frozen protocols reported here, neither
component recovers the payoff-based specialization criterion. High imitation
alone does not guarantee payoff-level specialization. Claims remain scoped to
the nominal $2$-vs-$2$, two-regime setting and training recipes studied here.

\end{document}
